\documentclass[8pt, aspectratio=169]{beamer}

% Theme
\usetheme[
  palette=darkred,                   % color profile: rose | darkred
  header=transparent,                % frame title: filled | transparent
  sidebar=left,
  progressbar=none,
  numbering=shifted,
  subsectionpage=progressbarextended,
]{crns}
\setbeamercovered{transparent}
\setbeamertemplate{section in toc}[sections numbered]

% Fonts: Kurier for text, Latin Modern for math
\usefonttheme{professionalfonts}     % keep beamer from overriding the math font
\usepackage{lmodern}
\usepackage{kurier}                  % after lmodern, so text stays in Kurier
\usepackage[french]{babel}

% Packages used by the showcase slides
\usepackage{amsmath, amssymb}
\usepackage{booktabs}

% Title page
\title{\huge Presentation Title}
\subtitle{Subtitle}
\author{\huge Achraf Ben-Hamadou}
\institute{\normalsize achraf.benhamadou@crns.rnrt.tn}
\date{\normalsize Assistant Professor\\Digital Research Centre of Sfax}

\begin{document}
\maketitle

\begin{frame}{Plan}
  \setbeamertemplate{section in toc}[sections numbered]
  \tableofcontents[hideallsubsections]
  %\tableofcontents
\end{frame}

\section{Section 1}
%--------------------------------------------------------------------------------------------------------------------------------
%--------------------------------------------------------------------------------------------------------------------------------
%--------------------------------------------------------------------------------------------------------------------------------
\begin{frame}{Slide's title}

\end{frame}
%--------------------------------------------------------------------------------------------------------------------------------
\begin{frame}[fragile]{Slide's title}{Slide's subtitle}
  A frame subtitle is given as a second argument to the frame:
  \begin{verbatim}\begin{frame}{Slide's title}{Slide's subtitle}\end{verbatim}
  or with \verb|\framesubtitle{...}| inside the frame.
\end{frame}
%--------------------------------------------------------------------------------------------------------------------------------
%--------------------------------------------------------------------------------------------------------------------------------
%--------------------------------------------------------------------------------------------------------------------------------
\section{Section 2}
\subsection{Sub-section 2.1}
%--------------------------------------------------------------------------------------------------------------------------------

\section{Elements}

\begin{frame}[fragile]{Typography}
      \begin{verbatim}The theme provides sensible defaults to
\emph{emphasize} text, \alert{accent} parts
or show \textbf{bold} results.\end{verbatim}

  \begin{center}becomes\end{center}

  The theme provides sensible defaults to \emph{emphasize} text,
  \alert{accent} parts or show \textbf{bold} results.
\end{frame}

\begin{frame}{Font feature test}
  \begin{itemize}
    \item Regular
    \item \textit{Italic}
    \item \textsc{Small Caps}
    \item \textbf{Bold}
    \item \textbf{\textit{Bold Italic}}
    \item \textbf{\textsc{Bold Small Caps}}
    \item \texttt{Monospace}
    \item \texttt{\textit{Monospace Italic}}
    \item \texttt{\textbf{Monospace Bold}}
    \item \texttt{\textbf{\textit{Monospace Bold Italic}}}
  \end{itemize}
\end{frame}

\begin{frame}{Mathematics}
  \begin{columns}[T,onlytextwidth]
    \column{0.5\textwidth}
      Fourier transform and Basel problem
      \[
        \hat{f}(\xi) = \int_{-\infty}^{+\infty} f(x)\, e^{-2\pi i x \xi}\, \mathrm{d}x
        \qquad
        \sum_{n=1}^{\infty} \frac{1}{n^2} = \frac{\pi^2}{6}
      \]
      Variational lower bound
      \begin{align*}
        \log p_\theta(\mathbf{x}) \geq{}&
        \mathbb{E}_{q_\phi(\mathbf{z} \mid \mathbf{x})}
        \big[ \log p_\theta(\mathbf{x} \mid \mathbf{z}) \big] \\
        &- D_{\mathrm{KL}}\big( q_\phi(\mathbf{z} \mid \mathbf{x}) \,\|\, p(\mathbf{z}) \big)
      \end{align*}
      Scaled dot-product attention
      \[
        \operatorname{Attention}(Q, K, V) =
        \operatorname{softmax}\!\left( \frac{Q K^{\top}}{\sqrt{d_k}} \right) V
      \]

    \column{0.5\textwidth}
      Diffusion forward process
      \begin{align*}
        q(\mathbf{x}_t \mid \mathbf{x}_{t-1}) &=
          \mathcal{N}\big(\mathbf{x}_t ; \sqrt{1-\beta_t}\, \mathbf{x}_{t-1}, \beta_t \mathbf{I}\big) \\
        q(\mathbf{x}_t \mid \mathbf{x}_0) &=
          \mathcal{N}\big(\mathbf{x}_t ; \sqrt{\bar{\alpha}_t}\, \mathbf{x}_0, (1-\bar{\alpha}_t) \mathbf{I}\big) \\
        \text{with} \quad \bar{\alpha}_t &= \prod_{s=1}^{t} (1-\beta_s)
      \end{align*}
      Matrices and cases
      \[
        \mathbf{R}_\theta =
        \begin{pmatrix}
          \cos\theta & -\sin\theta \\
          \sin\theta & \cos\theta
        \end{pmatrix}
        \quad
        \operatorname{ReLU}(x) =
        \begin{cases}
          x & \text{if } x > 0 \\
          0 & \text{otherwise}
        \end{cases}
      \]
      Alphabets
      \[
        \alpha\beta\gamma\delta\epsilon\lambda\mu\sigma\omega
        \quad \Gamma\Delta\Theta\Lambda\Sigma\Omega
        \quad \mathbb{R}\mathbb{N}\mathbb{Z}
        \quad \mathcal{L}\mathcal{N}\mathcal{X}
        \quad \mathbf{x}\mathbf{W}
      \]
  \end{columns}
\end{frame}

\begin{frame}{Lists}
  \begin{columns}[T,onlytextwidth]
    \column{0.33\textwidth}
      Items
      \begin{itemize}
        \item Milk \item Eggs \item Potatoes
      \end{itemize}

    \column{0.33\textwidth}
      Enumerations
      \begin{enumerate}
        \item First, \item Second and \item Last.
      \end{enumerate}

    \column{0.33\textwidth}
      Descriptions
      \begin{description}
        \item[PowerPoint] Meeh. \item[Beamer] Yeeeha.
      \end{description}
  \end{columns}
\end{frame}
\begin{frame}{Animation}
  \begin{itemize}[<+- | alert@+>]
    \item \alert<4>{This is\only<4>{ really} important}
    \item Now this
    \item And now this
  \end{itemize}
\end{frame}

\begin{frame}{Figures}
  \begin{figure}
    \newcounter{density}
    \setcounter{density}{20}
    \begin{tikzpicture}
      \def\couleur{alerted text.fg}
      \path[coordinate] (0,0)  coordinate(A)
                  ++( 90:5cm) coordinate(B)
                  ++(0:5cm) coordinate(C)
                  ++(-90:5cm) coordinate(D);
      \draw[fill=\couleur!\thedensity] (A) -- (B) -- (C) --(D) -- cycle;
      \foreach \x in {1,...,40}{%
          \pgfmathsetcounter{density}{\thedensity+20}
          \setcounter{density}{\thedensity}
          \path[coordinate] coordinate(X) at (A){};
          \path[coordinate] (A) -- (B) coordinate[pos=.10](A)
                              -- (C) coordinate[pos=.10](B)
                              -- (D) coordinate[pos=.10](C)
                              -- (X) coordinate[pos=.10](D);
          \draw[fill=\couleur!\thedensity] (A)--(B)--(C)-- (D) -- cycle;
      }
    \end{tikzpicture}
    \caption{Rotated square from
    \href{http://www.texample.net/tikz/examples/rotated-polygons/}{texample.net}.}
  \end{figure}
\end{frame}


\begin{frame}{Tables}
  \begin{table}
    \caption{Largest cities in the world (source: Wikipedia)}
    \begin{tabular}{@{} lr @{}}
      \toprule
      City & Population\\
      \midrule
      Mexico City & 20,116,842\\
      Shanghai & 19,210,000\\
      Peking & 15,796,450\\
      Istanbul & 14,160,467\\
      \bottomrule
    \end{tabular}
  \end{table}
\end{frame}
\begin{frame}{Blocks}
  Three different block environments are pre-defined and may be styled with an
  optional background color.

  \begin{columns}[T,onlytextwidth]
    \column{0.5\textwidth}
      \begin{block}{Default}
        Block content.
      \end{block}

      \begin{alertblock}{Alert}
        Block content.
      \end{alertblock}

      \begin{exampleblock}{Example}
        Block content.
      \end{exampleblock}

    \column{0.5\textwidth}

      \crnsset{block=fill}

      \begin{block}{Default}
        Block content.
      \end{block}

      \begin{alertblock}{Alert}
        Block content.
      \end{alertblock}

      \begin{exampleblock}{Example}
        Block content.
      \end{exampleblock}
  \end{columns}
\end{frame}


\begin{frame}[standout, noframenumbering]
\huge{Thanks for your attention}
\end{frame}

\end{document}
%--------------------------------------------------------------------------------------------------------------------------------

